\begingroup
\renewcommand{\arraystretch}{1.45}
\fancyhead[L]{\small ARITHMETIC–GEOMETRIC MEAN}
\Needspace{32\baselineskip}
\AGMArticle{3}

To show more clearly how rapidly sequences (I) and (II) approach their limit
to the right, and how the former increases and the latter decreases to the left,
we give a few examples here:

\Needspace{20\baselineskip}\begin{center}\textit{Example} 1.\quad $a=1,\ b=0{,}2$\end{center}
\begingroup\small\setlength{\arraycolsep}{0pt}
\[\begin{array}{r@{\,=\,}l@{\quad}r@{\,=\,}l}
{}^{\prime\prime\prime\prime}a&15{,}83795\,47919\,02&{}^{\prime\prime\prime\prime}b&0{,}00000\,00000\,00000\,00005\,7\\[6pt]
{}^{\prime\prime\prime}a&7{,}91897\,73959\,512&{}^{\prime\prime\prime}b&0{,}00000\,00013\,481\\[6pt]
{}^{\prime\prime}a&3{,}95948\,86986\,4971&{}^{\prime\prime}b&0{,}00010\,30955\,7682\\[6pt]
{}^{\prime}a&1{,}97979\,58971\,13271\,23927\,9&{}^{\prime}b&0{,}02020\,41028\,86728\,76072\,1\\[6pt]
a&1{,}00000\,00000\,00000\,00000\,0&b&0{,}20000\,00000\,00000\,00000\,0\\[6pt]
a'&0{,}60000\,00000\,00000\,00000\,0&b'&0{,}44721\,35954\,99957\,93928\,2\\[6pt]
a''&0{,}52360\,67977\,49978\,96964\,1&b''&0{,}51800\,40128\,22268\,36005\,0\\[6pt]
a'''&0{,}52080\,54052\,86123\,66484\,5&b'''&0{,}52079\,78709\,39876\,24344\,0\\[6pt]
a''''&0{,}52080\,16381\,12999\,95414\,3&b''''&0{,}52080\,16380\,99375\\[6pt]
a^{\mathrm V}&0{,}52080\,16381\,06187&b^{\mathrm V}&0{,}52080\,16381\,06187
\end{array}
\]\endgroup

Here $a^{\mathrm V}$, $b^{\mathrm V}$ first differ in the $23^{\mathrm{rd}}$ digit;
${}^{\mathrm V}b$ is smaller than $(\frac{1}{10})^{40}$;
${}^{\prime\prime\prime\prime}a$, ${}^{\mathrm V}a$, ${}^{\mathrm{VI}}a$, etc.
form, for practical purposes, a geometric sequence whose common ratio is $=2$.

\Needspace{20\baselineskip}\begin{center}\textit{Example} 2.\quad $a=1,\ b=0{,}6$.\end{center}
\begingroup\small\setlength{\arraycolsep}{0pt}
\[\begin{array}{r@{\,=\,}l@{\quad}r@{\,=\,}l}
{}^{\prime\prime\prime\prime}a&14{,}35538\,2913&{}^{\prime\prime\prime\prime}b&0{,}00000\,00000\\[6pt]
{}^{\prime\prime\prime}a&7{,}17769\,14569\,307&{}^{\prime\prime\prime}b&0{,}00001\,73070\\[6pt]
{}^{\prime\prime}a&3{,}58885\,43819\,99831\,75712\,7&{}^{\prime\prime}b&0{,}01114\,56180\,00168\,24287\,3\\[6pt]
{}^{\prime}a&1{,}80000\,00000\,00000\,00000\,0&{}^{\prime}b&0{,}20000\,00000\,00000\,00000\,0\\[6pt]
a&1{,}00000\,00000\,00000\,00000\,0&b&0{,}60000\,00000\,00000\,00000\,0\\[6pt]
a'&0{,}80000\,00000\,00000\,00000\,0&b'&0{,}77459\,66692\,41483\,37703\,6\\[6pt]
a''&0{,}78729\,83346\,20741\,68851\,8&b''&0{,}78719\,58685\,06172\,16741\,6\\[6pt]
a'''&0{,}78724\,71015\,63456\,92796\,7&b'''&0{,}78724\,70999\\[6pt]
a''''&0{,}78724\,71007\,8&b''''&0.78724\,71007\,8
\end{array}
\]\endgroup
\SourceLeaf{373}
\Needspace{20\baselineskip}\begin{center}\textit{Example} 3.\quad $a=1,\ b=0{,}8$.\end{center}
\begingroup\small\setlength{\arraycolsep}{0pt}
\[\begin{array}{r@{\,=\,}l@{\quad}r@{\,=\,}l}
{}^{\mathrm V}a&25{,}19190\,722&{}^{\mathrm V}b&0{,}00000\,00000\,0\\[6pt]
{}^{\prime\prime\prime\prime}a&12{,}59595\,36116\,78&{}^{\prime\prime\prime\prime}b&0{,}00000\,00133\,367\\[6pt]
{}^{\prime\prime\prime}a&6{,}29797\,68125\,07655\,42373\,4&{}^{\prime\prime\prime}b&0{,}00040\,98644\,58278\,08440\,9\\[6pt]
{}^{\prime\prime}a&3{,}14919\,33384\,82966\,75407\,2&{}^{\prime\prime}b&0{,}05080\,66615\,17033\,24592\,8\\[6pt]
{}^{\prime}a&1{,}60000\,00000\,00000\,00000\,0&{}^{\prime}b&0{,}40000\,00000\,00000\,00000\,0\\[6pt]
a&1{,}00000\,00000\,00000\,00000\,0&b&0{,}80000\,00000\,00000\,00000\,0\\[6pt]
a'&0{,}90000\,00000\,00000\,00000\,0&b'&0{,}89442\,71909\,99915\,87856\,4\\[6pt]
a''&0{,}89721\,35954\,99957\,93928\,2&b''&0{,}89720\,92687\,32734\\[6pt]
a'''&0{,}89721\,14321\,16346&b'''&0{,}89721\,14321\,13738\\[6pt]
a''''&0{,}89721\,14321\,15042&b''''&0{,}89721\,14321\,15042
\end{array}
\]\endgroup

\Needspace{20\baselineskip}\begin{center}\textit{Example} 4.\quad $a=\surd 2,\ b=1$.\end{center}
\begingroup\small\setlength{\arraycolsep}{0pt}
\[\begin{array}{r@{\,=\,}l@{\quad}r@{\,=\,}l}
{}^{\prime\prime\prime\prime}a&19{,}17024\,37557\,69475\,31905\,0&{}^{\prime\prime\prime\prime}b&0{,}00000\,00009\,32560\,02627\,6\\[6pt]
{}^{\prime\prime\prime}a&9{,}58512\,18783\,51017\,67266\,3&{}^{\prime\prime\prime}b&0{,}00013\,37064\,06056\,69181\,0\\[6pt]
{}^{\prime\prime}a&4{,}79262\,77923\,78537\,18223\,7&{}^{\prime\prime}b&0{,}03579\,93323\,67652\,95745\,7\\[6pt]
{}^{\prime}a&2{,}41421\,35623\,73095\,04880\,2&{}^{\prime}b&0{,}41421\,35623\,73095\,04880\,2\\[6pt]
a&1{,}41421\,35623\,73095\,04880\,2&b&1{,}00000\,00000\,00000\,00000\,0\\[6pt]
a'&1{,}20710\,67811\,86547\,52440\,1&b'&1{,}18920\,71150\,02721\,06671\,7\\[6pt]
a''&1{,}19815\,69480\,94634\,29555\,9&b''&1{,}19812\,35214\,93120\,12260\,7\\[6pt]
a'''&1{,}19814\,02347\,93877\,20908\,3&b'''&1{,}19814\,02346\,77307\,20579\,8\\[6pt]
a''''&1{,}19814\,02347\,35592\,20744\,1&b''''&1{,}19814\,02347\,35592\,20743\,9
\end{array}
\]\endgroup

\Needspace{35\baselineskip}
\AGMArticle{4}

In addition to the sequences
\[
\begin{array}{l}
.\;.\;.{}^{\prime\prime\prime}a,\;{}^{\prime\prime}a,\;{}^{\prime}a,\; a,\; a',\; a'',\; a'''.\;.\;.\\[6pt]
.\;.\;.{}^{\prime\prime\prime}b,\;{}^{\prime\prime}b,\;{}^{\prime}b,\; b,\; b',\; b'',\; b'''.\;.\;.
\end{array}
\]
let there be two others,
\[
\begin{array}{l}
.\;.\;.{}^{\prime\prime\prime}c,\;{}^{\prime\prime}c,\;{}^{\prime}c,\; c,\; c',\; c'',\; c'''.\;.\;.\\[6pt]
.\;.\;.{}^{\prime\prime\prime}d,\;{}^{\prime\prime}d,\;{}^{\prime}d,\; d,\; d',\; d'',\; d'''.\;.\;.
\end{array}
\]
formed in the same way; and let us suppose that two terms of the latter are
\textit{proportional} to two terms of the former, for example, $a:b=c:d$, or $a:c=b:d=1:n$.
Then all the terms of I will stand to the terms of III, and all those of II to those of IV
(taken in corresponding positions relative to $a$, $b$, $c$, $d$) in the same ratio; thus
\SourceLeaf{374}
\[
\begin{gathered}c'=na',\quad c''=na'',\quad c'''=na''',\\ {}^{\prime}c={}^{\prime}an,\quad d^{\mathrm V}=nb^{\mathrm V},\quad{}^{\mathrm V}d={}^{\mathrm V}bn\end{gathered}
\]
It follows easily that the limit of sequences I, II will also stand to the limit of sequences
III, IV as $1$ to $n$, or, in general, $\mathrm M(na,nb)=n\mathrm M(a,b)$.
Consequently, in general, $\mathrm M(a,b)=a\mathrm M(1,\frac{b}{a})=b\mathrm M(\frac{a}{b},1)$.

\clearpage
\endgroup
